%===== Chapter 1: Project Framework and Context =====%

\chapter{Project Framework and Context}

\section*{Introduction}
This chapter introduces the host company, presents the project, analyses the
existing market solutions, and justifies the chosen development methodology and
technology stack.

\section{Host Company: Anter Lab}

\textbf{Anter Lab} is a Tunisian technology company founded by two entrepreneurs,
headquartered in Montplaisir, Tunis. The company empowers Small and Medium Enterprises
(SMEs) with smart HR solutions and delivers web, mobile, and marketing services,
blending innovation with creative strategy.

The company's dual focus on both product development and service delivery creates a
fertile environment for ambitious engineering projects such as Malltech.

\begin{longtable}{|m{3.5cm}|m{10cm}|}
    \hline
    \textbf{Field} & \textbf{Details} \\
    \hline
    \endhead
    \endfoot
    \endlastfoot
    \hline
    Company & Anter Lab \\
    \hline
    Type & Private technology company \\
    \hline
    Founded & Tunis, Tunisia \\
    \hline
    Address & Montplaisir, Bab El Bhar, Tunis, 1073 \\
    \hline
    Services & HR SaaS, web/mobile development, digital marketing \\
    \hline
    \caption{Anter Lab company profile}
    \label{tab:company_profile}
\end{longtable}

\section{Project Presentation}

\textbf{Malltech} is a multi-service e-commerce platform designed for the Tunisian
market. The name combines \textit{mall} (a unified commercial space) with
\textit{tech} (modern software engineering). The platform targets two audiences
simultaneously:

\begin{itemize}
    \item \textbf{End customers} — who browse products, manage a cart and wishlist, place orders, and track deliveries through the storefront.
    \item \textbf{Merchants and administrators} — who manage the entire catalog, order lifecycle, customer base, promotional campaigns, content, and business configuration through the backoffice dashboard.
\end{itemize}

\section{Study of Existing Solutions}

Before designing Malltech, a comparative analysis of available platforms was conducted:

\begin{longtable}{|m{2.8cm}|m{1.5cm}|m{4cm}|m{4.5cm}|}
    \hline
    \textbf{Platform} & \textbf{Type} & \textbf{Strengths} & \textbf{Weaknesses for Tunisia} \\
    \hline
    \endhead
    \endfoot
    \endlastfoot
    \hline
    Shopify & SaaS & Polished UX, large ecosystem & Monthly cost in USD, no local payment context \\
    \hline
    WooCommerce & Open-source & Free, flexible & Requires WordPress, hard to scale \\
    \hline
    Magento & Open-source & Enterprise-grade & Very heavy, expensive hosting \\
    \hline
    Custom PHP & Bespoke & Tailored & No modern stack, no CI/CD \\
    \hline
    Social media & Free & Zero setup & No inventory, no analytics \\
    \hline
    \caption{Comparison of existing e-commerce solutions}
    \label{tab:existing_solutions}
\end{longtable}

\textbf{Conclusion:} None of the existing solutions offered a modern TypeScript
full-stack architecture, native support for Tunisian commercial workflows, and the
ability to be self-hosted affordably. Building a custom platform was justified.

\section{Development Methodology}

\subsection{Chosen Approach: Iterative and Incremental}

Given that this is a solo development project with evolving requirements, a
\textbf{lightweight iterative approach} was adopted, inspired by Agile principles
but without formal sprint ceremonies. Development was organised into functional
milestones:

\begin{enumerate}
    \item \textbf{Milestone 1} — Core infrastructure: project scaffolding, Docker setup, authentication, user models.
    \item \textbf{Milestone 2} — Product catalog: categories, subcategories, brands, products, stock management, 3D model support.
    \item \textbf{Milestone 3} — Commerce engine: basket, orders, promotions, coupons, shipping rates, wishlist.
    \item \textbf{Milestone 4} — Financial layer: invoices, credit notes, PDF generation, selling points.
    \item \textbf{Milestone 5} — Customer storefront: all public-facing Next.js pages.
    \item \textbf{Milestone 6} — Admin dashboard: all backoffice Next.js pages.
    \item \textbf{Milestone 7} — Ancillary services: support bot, email notifications, CMS, analytics.
    \item \textbf{Milestone 8} — Infrastructure: CI/CD pipelines, monitoring, auto-healing, Caddy reverse proxy.
\end{enumerate}

\subsection{Version Control}

Git was used for version control across three GitHub repositories (one per service).
Each repository has its own independent CI/CD pipeline via GitHub Actions.

\section{Technology Choices Rationale}

\subsection{Backend — NestJS + MongoDB}

\textbf{NestJS} was chosen for its opinionated module system, built-in dependency
injection, decorator-based route definitions, and first-class TypeScript support.
These features promote code organisation and maintainability in a large codebase.

\textbf{MongoDB} was chosen because the product catalog (nested attributes, variable
schema per category, embedded stock) maps naturally to documents. Mongoose provides
schema validation at the application layer without requiring migrations.

\subsection{Frontend and Admin — Next.js}

\textbf{Next.js 16} with the App Router was chosen for server-side rendering (SSR)
and static generation (SSG) for SEO, React Server Components for reduced client-side
JavaScript, built-in image optimisation, and a unified framework for both the
storefront and the admin.

\subsection{Infrastructure — Docker + Caddy + GitHub Actions}

\textbf{Docker} isolates each service and guarantees environment parity between
local development and production. \textbf{Caddy} provides zero-configuration
automatic HTTPS via Let's Encrypt. \textbf{GitHub Actions} delivers CI/CD directly
integrated with the source repositories at no additional tooling cost.

\section{Project Architecture Overview}

Malltech is composed of three independently deployable services, all running inside
a shared Docker network on a single VPS:

\begin{figure}[H]
\centering
% Replace with actual architecture diagram when available
\fbox{\parbox{0.85\textwidth}{\centering\vspace{2cm}
\textit{[Architecture Diagram Placeholder]}\\
\small Internet $\rightarrow$ Caddy (80/443) $\rightarrow$ \{Frontend :3000, Admin :3000, API :3000\}\\
API $\rightarrow$ MongoDB\\
Netdata (monitoring)
\vspace{2cm}}}
\caption{High-level deployment architecture of Malltech}
\label{fig:architecture}
\end{figure}

\section*{Conclusion}
This chapter presented the host company Anter Lab, the motivation behind Malltech,
justified the technology choices through a competitive analysis, and described the
iterative development methodology used throughout the project. The following chapter
details the functional and non-functional requirements identified for the platform.
